% Guia do professor — Capítulo 3: Termodinâmica da Atmosfera
\section{Capítulo 3 --- Termodinâmica da Atmosfera}

\subsection*{Visão geral e ritmo}

O capítulo transforma a 1ª lei da termodinâmica em quatro ferramentas:
gradiente adiabático seco, temperatura potencial, balanço euleriano de
calor e balanço de superfície. \textbf{Ritmo sugerido: 2 aulas de 2\,h.}

\begin{itemize}
  \item \textbf{Aula 1} — calores específicos e latente, 1ª lei na forma
        meteorológica, $\Gamma_d$, temperatura potencial (com o exemplo
        föhn). Fechar com o Caso~1 (10\,min).
  \item \textbf{Aula 2} — derivada material e advecção, balanço euleriano,
        balanço de superfície e Bowen. Fechar com Casos~2--3 (15--20\,min;
        o Caso~3 merece a projeção mais longa).
\end{itemize}

\subsection*{Objetivos de aprendizagem}

\begin{enumerate}
  \item aplicar $\delta q = c_p dT - \alpha dP$ e explicar fisicamente cada
        termo;
  \item deduzir $\Gamma_d = g/c_p$ e distingui-lo do gradiente ambiental;
  \item calcular $\theta$, justificar sua conservação e usá-la para comparar
        parcelas de níveis diferentes;
  \item escrever o balanço euleriano, estimar advecção com números reais e
        entender a derivada material;
  \item repartir $F^*$ em $F_H$, $F_E$, $F_G$ via razão de Bowen e prever o
        caráter do ciclo diurno de uma superfície.
\end{enumerate}

\subsection*{Pré-requisitos a reativar}

1ª lei na forma $dU = \delta Q - PdV$ (física básica); hidrostática e gás
ideal do Cap.~1 (usados nas duas deduções centrais); regra da cadeia em
várias variáveis (derivada material). Alunos de física têm tudo — o novo é
o \emph{vocabulário} meteorológico.

\subsection*{Armadilhas conceituais frequentes}

\begin{itemize}
  \item \textbf{``Esfria porque lá em cima é frio''}: a parcela esfria por
        \emph{expansão}, não por contato. É a confusão número 1 do capítulo.
  \item \textbf{$\Gamma_d$ vs.\ gradiente ambiental}: 9{,}8 é o que a
        parcela faz; 6{,}5 é a média do que a sondagem mostra. Confundi-los
        inviabiliza o Cap.~5.
  \item \textbf{$\theta$ como ``temperatura exótica''}: insistir que é só um
        rótulo conservado — comparar com energia mecânica em queda livre.
  \item \textbf{Sinal da advecção}: vento soprando \emph{do} quente
        \emph{para} o frio dá $-U\,\partial T/\partial x > 0$ (aquece).
        Fazer o desenho do gradiente com setas antes da fórmula.
  \item \textbf{Bowen como propriedade fixa}: $\beta$ depende da umidade
        disponível — o mesmo solo seco ou molhado muda $\beta$ por ordem de
        magnitude.
\end{itemize}

\subsection*{Demonstração de baixo custo}

Borrifador de álcool no termômetro (ou no dorso da mão): queda imediata de
vários graus — calor latente em ação. Alternativa: bomba manual de encher
pneu aquece ao comprimir (adiabática no sentido oposto).

\subsection*{Problemas sugeridos do livro (seção de exercícios do Cap.~3)}

Aquecimento: conversões de calor/temperatura, cálculo de $\theta$ para
vários níveis, advecção com vento e gradiente dados. Aprofundamento:
problemas de balanço de superfície com Bowen; síntese: ``o que aconteceria
se $L_v$ da água fosse 10× menor?''. \emph{Confira a numeração na sua cópia
(v1.02b) antes de publicar a lista.}

\subsection*{Uso do notebook \texttt{cap03\_termodinamica.ipynb}}

\begin{center}
\begin{tabular}{lll}
\hline
Caso & Conteúdo & Momento sugerido \\
\hline
1 & Parcela adiabática (EDO); conservação de $\theta$ & fim da Aula 1 \\
2 & Ciclo diurno: balanço de superfície com Bowen & Aula 2, após Bowen \\
3 & Advecção 1D (EDP, esquema \emph{upwind}) & Aula 2, após advecção \\
4 & $\theta$ da sondagem real; camada de mistura & Aula 2 ou tarefa \\
\hline
\end{tabular}
\end{center}

O Caso~3 é o momento didático mais rico: é a primeira EDP integrada no
curso (``previsão de verdade''), e a difusão numérica visível abre a
conversa sobre por que modelos operacionais usam esquemas sofisticados
(Cap.~20).

\subsection*{Exercícios complementares e soluções}

Nas notas: T1--T5 ($\Gamma_d$ terrestre e marciano, dedução de $\theta$,
$\theta$--entropia, solução geral da advecção, partição de Bowen) e N1--N4
(células reservadas no notebook). Soluções em
\texttt{cap03\_solucoes.ipynb} (\emph{não distribuir}): T2--T4 com
verificação simbólica \texttt{sympy}; N3 quantifica a difusão numérica do
\emph{upwind}. Pares teoria--numérica: T2$\leftrightarrow$N1 e
T4$\leftrightarrow$N3.

\subsection*{Conexões com o resto do curso}

$\Gamma_d$ e $\theta$ $\to$ estabilidade e diagramas termodinâmicos
(Cap.~5); calor latente $\to$ umidade (Cap.~4) e tempestades (Cap.~14);
advecção $\to$ frentes (Cap.~12) e NWP (Cap.~20); balanço de superfície
$\to$ camada limite (Cap.~18).
